\paragraph{Finite denominators at abelian reduction.}
\label{intjet:sec}
Write \(d=\deg P=\sum a_{\tau i}\), \(\Delta=|\operatorname{disc}E|\),
and \(v\) for the additive valuation of a mixed-characteristic
discrete valuation ring \(R\). Valuations and all the inclusions below
may be tested after finite extension. We retain the integral symmetric
line, its rigidification, the section, the Hodge directions and the
polarized biextension on an abelian scheme \(A/R\) as in
Proposition~\ref{newtheta:goodjet}. The graph of \(w_0\in O_E\)
in that biextension gives the line. In a semiabelian degeneration the
generic abelian scheme need not extend in this form.

\begin{lemma}[Integral Taylor derivatives of harmonic jets]\label{intjet:der}
In the preceding formal chart take an integral raw coefficient,
any nonnegative tensor power of the line, and an integral homogeneous
harmonic polynomial of degree \(d\) at the standard lattice.
Its evaluated jet belongs to \(\Delta^{-\lfloor d/2\rfloor}R\).
The assertion holds in every residue characteristic.
\end{lemma}
\begin{proof}
Work over the fraction field in logarithm coordinates. For a series
\(F\) define its Taylor values to be the mixed derivatives
\(\partial^\alpha F(0)\), with no division by \(\alpha!\).
If \(F\) comes from an integral function on the formal group,
on integral tangent vectors these values lie in \(R\). Indeed
each vector extends by translation to an invariant integral
derivation of the formal group ring. In logarithm coordinates it
is the constant vector field, so applying a sequence of these
derivations and evaluating at zero gives exactly the desired value.
This argument works on a product and on any list of integral
directions, whether or not they form an integral basis.

Series with integral Taylor values form a ring by Leibniz's formula.
An inverse with unit constant term has integral Taylor values by the
same formula and induction. For \(g(0)=1\) and integral Taylor
values, those of \(\log g\) are integral by
\(\partial_i\log g=g^{-1}\partial_i g\). For \(h(0)=0\) with
integral Taylor values, use
\(\partial_i\exp h=(\exp h)\partial_i h\) and induction.
These are identities over the fraction field and give integrality
of derivatives without an estimate for any inverse factorial.

Let \(h(X,Y)\) be the logarithmic correction of the biextension in
the chosen frame. Its two additive defects and its \(O_E\)-adjunction
defect are logarithms of integral units with constant term one, as
in Proposition~\ref{newtheta:goodjet}. For an index \(\alpha\) of
degree at least two write \(\alpha=\alpha'+e_j\), \(|\alpha'|>0\).
Apply \(\partial_X^{\alpha'}\partial_{X'_j}\partial_Y^\beta\) to
\(h(X,Y)+h(X',Y)-h(X+X',Y)\) at zero. The result is
\(-\partial_X^\alpha\partial_Y^\beta h(0)\). This proves its
integrality. Interchanging the two arguments treats \(|\beta|\ge2\).
The two axes vanish. In bidegree \((1,1)\) write \(c_{\tau i,\nu j}\)
for the derivative along eigenvectors of types \(\tau,\nu\).
The adjunction defect gives
\((\tau(a)-\nu(\bar a))c_{\tau i,\nu j}\in R\) for every \(a\in O_E\).
If \(\tau\ne\bar\nu\), then \(\Delta c_{\tau i,\nu j}\in R\).
Indeed choose an integral basis \(b_i\) and its trace-dual basis
\(b_i^\vee\). The embedding matrices are inverse transposes, so
\(\sum_i\tau(b_i^\vee)\sigma(b_i)\) is one for \(\sigma=\tau\)
and zero otherwise. All \(\Delta b_i^\vee\) are integral by the
adjugate of the trace matrix. Apply this equality to
\(\sigma=\tau,\bar\nu\) and subtract.

Thus the gauge on the eigen-directions is the sum of a series with
integral Taylor values, a bilinear form with coefficients in
\(\Delta^{-1}R\) on the nonconjugate pairs, and an unrestricted
bilinear form on the conjugate pairs. On the graph of \(w_0\) these
are quadratic terms, summed over the product factors. The last
quadratics are killed by harmonicity as in
Proposition~\ref{newtheta:goodjet}. In the exponential of the
nonconjugate quadratic each derivative of order \(j\) at zero is
a sum over partitions into pairs of the \(j\) labelled derivative
positions. This follows by differentiating \(\partial_i\exp
Q=(\exp Q)\partial_i Q\); a singleton has zero value and derivatives
of \(Q\) of order greater than two vanish. Each summand costs at
most \(\lfloor j/2\rfloor v(\Delta)\). Leibniz's formula for the
remaining integral factors gives the claimed bound.
\end{proof}

\begin{proposition}[All abelian tests at a locally standard lattice]
\label{intjet:standard}
Suppose \(\Lambda_\ell=(O_E^n)_\ell\) and the integral abelian data
above are given. For every allowed \(N\) and each pair of translation
points in the original order one has
\[
 v(J_{P,N}(u,z))\ \ge\
 d\bigl(v(N)-v(M)\bigr)-\lfloor d/2\rfloor v(\Delta).
\]
In particular \(D_M^2\) clears the rescaled sum at every such test,
including the level primes, \(2\), and primes ramified in \(E\).
\end{proposition}
\begin{proof}
Let \(N_0=\operatorname{lcm}(M,\exp(O_E^n/\Lambda_0))\) as before.
Now \(v(N_0)=v(M)\). The identity
\(J_{P,N}= (N/N_0)^d J_{P,N_0}\) on the generic fibre proved in
Proposition~\ref{newtheta:goodjet} works whether or not \(\ell\mid N_0\).
All torsion points of the smaller test extend after finite extension
by properness. For \(t\in A[N_0^2]\) the line \(T_t^*L^{N_0^2}\)
is integrally isomorphic to \(L^{N_0^2}\). The class of their
ratio in the dual abelian scheme is zero since \(N_0^2t=0\),
and an invertible module on the local base is free. Any integral
translation lift has its \(N_0^2\)-th power in \(R^\times\).
The generic scalar adjusting it to the normalized lift therefore
has valuation zero. This applies to both translations in succession,
also for torsion of residue-characteristic-power order.
The smaller-level section has integral raw germ. Apply
Lemma~\ref{intjet:der}. Local standardness makes \(P\) integral in
the standard columns and \(s\) integral at this place. Sum after the
index cancellation of Lemma~\ref{newtheta:index} and use
\(D_M=C_M(M\Delta)^d\).
\end{proof}

\begin{theorem}[The trace at an arbitrary self-dual lattice at abelian reduction]
\label{intjet:trace}
Suppose the abelian data above extend integrally. Take any
\(\Lambda,N,P,s\) in the rescaled series, with \(\Lambda\) standard
at the primes of \(M\), and the prescribed level points \(t_u\).
At a prime \(\ell\nmid M\),
\[
 v\left(s\sum_{\bar z\in Z_\Lambda/K_0}J_{P,N}(u,\bar z)\right)
 \ \ge\ d\,v(N)-\lfloor d/2\rfloor v(\Delta)
 \qquad(u\in\mathcal Q).
\]
Together with Proposition~\ref{intjet:standard}, this proves the
finite denominator test at every residue characteristic on these
abelian models.
\end{theorem}
\begin{proof}
Let \(O=O_E\), \(B=A^n\), \(B_0=\Lambda_0\otimes_O A\) and
\(B_1=\Lambda\otimes_O A\). A full lattice over the Dedekind
ring \(O\) is finitely generated projective. Its tensor with \(A\)
is obtained by splitting an idempotent on a product of copies of
\(A\). This gives abelian schemes over \(R\). Their Lie algebras
are obtained by taking the same tensor with \(\operatorname{Lie}A\).
Consider the maps induced by
\[
 O^n\xrightarrow{\,N\,}\Lambda_0
 \xrightarrow{\,i\,}\Lambda,\qquad
 \Lambda_0\xrightarrow{\,j\,}O^n
\]
where \(i,j\) are inclusions. Write \(a_0,i,j\) also for the
resulting isogenies, and \(a=i a_0\). Then \(j a_0=[N]\).
These maps are isogenies on every fibre since suitable integer
multiplications factor through them in both directions. They
are finite flat by properness and flatness of a quasi-finite map
between smooth schemes of the same relative dimension.

Give each tensor scheme its rigidified line from the graph of
\(w_0\) in the polarized biextension for the restricted Hermitian
form. Denote the lines on \(B,B_1\) by \(\mathcal L,\mathcal L_1\).
Then \(j^*\mathcal L=i^*\mathcal L_1\) and
\(a^*\mathcal L_1=[N]^*\mathcal L=\mathcal L^{N^2}\)
by rigidified identifications. To see the normalization, on free
modules a matrix \(U\) pulls the biextension for a Hermitian
matrix \(G\) back to that for \(U^\dagger G U\). Expand the
biextension in its two arguments and move each first-argument
entry across by Rosati adjunction. These identifications are
the identity on the axes. Restriction to direct summands proves
the assertion for projective modules as well. This construction
uses the line on the graph, so no choice among symmetric lines
inducing the same polarization is needed. See Amir-Khosravi,
\S\S1.2--1.4~\cite{serreTensor}.

On the geometric generic fibre, the period lattice calculation of
Lemma~\ref{newtheta:index} gives \(\ker a=Z_\Lambda\) and
\(\ker a_0=K_0\). The natural action of \(\ker a\) on the pullback
line is exactly its normalized lift. It has \(N^2\)-th power one,
and inversion on the symmetric pullback sends the lift to its
inverse. These conditions determine the lift since \(N\) is odd.
In particular they agree with the products of the individual
normalized lifts on the \(n\) factors by the same uniqueness.
Put \(t_0=j^*(s_0^{\boxtimes n})\), viewed as a section of
\(i^*\mathcal L_1\). Before the level translation the sum of
the \(|Z_\Lambda/K_0|\) translated sections on \(B\) is
\[
                 a^*\operatorname{Tr}_i(t_0).
\]
Indeed a finite étale isogeny on the generic fibre has trace
equal, after pullback, to the sum over its kernel. The input
on \(B\) was \(a_0^*t_0=[N]^*(s_0^{\boxtimes n})\).

The determinant of \(di\) generates the ideal
\(\mathfrak b R=\prod_\tau
 (\tau(\det\Lambda_0\,\det\Lambda^{-1})R)^{m_\tau}\).
Choose bases of the two lattices over the semilocal Dedekind
ring \(O\otimes\mathbf Z_{(\ell)}\), where they are free.
The determinant on each generic eigenspace is the indicated
embedding of the determinant of the inclusion, repeated \(m_\tau\)
times. The chosen bases give integral changes of basis on the
Lie lattices, so this calculates their determinant ideal even
when the eigenspaces do not split integrally.
Translation of invariant differential frames carries this same
different ideal to every point. For a finite flat map between smooth
schemes of equal relative dimension the different is generated by
the determinant of its differential. Locally in the finite
algebra the trace functional is this generator times a generator
of the dualizing module, by the complete intersection formula
for the different~\cite{tateDifferent}. Trivializing \(\mathcal L_1\)
shows that \(\operatorname{Tr}_i(t_0)\) is divisible by
\(\mathfrak b R\) as a section. Thus \(S=s\operatorname{Tr}_i(t_0)\)
is integral, including when \(v(s)<0\).

It remains to apply the actual final translation. Its point
\(t=t_u\) is killed by a power of the half-level ideal, hence by
a power of \(M\); this extra condition is part of the evaluation
rule. In particular its order \(h\) is prime to \(\ell\).
\src{src/Formalism/LocalToGlobal/ProjectiveLattices/ThShort.lean}{252}{The support of the prescribed division point}.
Set \(b=a(t)\). Rigidify the line
\(\mathcal D=T_b^*\mathcal L_1\otimes\mathcal L_1^{-1}\) at zero
by choosing an integral generator of its fibre. It is an element
of \(\operatorname{Pic}^0(B_1)[h]\), since \(b\mapsto
T_b^*\mathcal L_1\otimes\mathcal L_1^{-1}\) is a homomorphism
to the dual abelian scheme. Thus \(\mathcal D^h\) has an integral
rigidified trivialization. On the formal neighborhood of zero
there is a unique \(h\)-th root frame with prescribed value at
zero and this trivialization as its \(h\)-th power. Coefficient
recursion only divides by the unit \(h\).

The normalized lift at \(t\) identifies the translation of
\(a^*S\) with the pullback of \(T_b^*S\), followed by an
isomorphism \(a^*\mathcal D\to\mathcal O_B\). This isomorphism
is integral by the translation-lift argument of
Proposition~\ref{intjet:standard}. Its \(h\)-th power and the
pullback of the chosen trivialization differ by a global unit
of the abelian scheme, hence by an element of \(R^\times\).
It takes the pullback of the root frame to a constant unit:
over the fraction field a formal unit whose \(h\)-th power is
constant is constant. In particular, in the formal neighborhood
the translated section is, up to a constant unit, \(a^*\) of
an integral germ of \(\mathcal L_1\), obtained by dividing the
germ \(T_b^*S\) by the root frame. This explains why the final
translation does not change the trace divisibility.

To compute the jet of this pullback, choose a split surjection
\(U:O^t\to\Lambda\) with section \(V:\Lambda\to O^t\),
extended rationally, and put \(G=U^\dagger U\). Its entries
are integral by self-duality. View \(U\) as an \(n\)-by-\(t\)
matrix in \(E^n\). Then \(UV=I_n\), \(V^\dagger G V=I_n\),
and \(NV\) has integral entries. The map \(a\) factors as \(NV\)
followed by \(U\). Pull the integral germ back by \(U\); its raw
coefficient in an integral frame of the line for \(G\) is integral.
In the corrected frame the chain rule for \(NV\) gives \(N^d\)
times the differential operator \(P_V\) obtained from \(P\)
by sending a column-direction \(v\) of type \(\tau\) to
\(\tau(V)v\). Indeed in logarithm coordinates the map is
linear on each eigenspace. The corrected biextension frame
multiplies in both arguments and intertwines the adjoint
\(O\)-actions by its defining identities. Expanding the matrix
entries in these identities identifies its pullback with the
corrected frame for \((NV)^\dagger G(NV)=N^2 I_n\).
Thus the comparison differentiates the full section with the
same correction as in the original jet.

The coefficients of \(P_V\) are integral. The defining integral
polynomial ring is generated by the linear directions from
\(\Lambda\), and \(V\Lambda\subset O^t\).
\src{src/Formalism/LocalToGlobal/ProjectiveLattices/ThShort.lean}{235}{Integral polynomial directions of a nonstandard lattice}.
For the Hermitian matrix \(G\) the quadratic ambiguity in the
corrected log frame comes from summing
\(h(X_j,w_0G_{jk}X_k)\) over its columns. For a conjugate pair
of types its contraction on \(P_V\) uses the coefficients
\(\bar\tau(G_{jk})\). By the ordinary chain rule and
\(V^\dagger G V=I_n\), it is the image under \(V\) of the standard
contraction on \(P\); hence it vanishes. All other terms obey
Lemma~\ref{intjet:der} because the entries of \(G\) and \(w_0\)
are integral. This costs at most \(\lfloor d/2\rfloor v(\Delta)\).
The factor \(N^d\) proves the stated inequality.
\end{proof}

\paragraph{The norm of the generator at semiabelian reduction.}
Let \(\mathcal A/R\) be the semiabelian model in an integral test,
with abelian generic fibre \(A\), rigidified cubical line
\(\mathcal L\), integral section \(s_0\), and integral tangent
vectors. Its principal polarization is \(\lambda\). Write
\(\mathcal P\) for the rigidified Poincaré line of \(A\times A^t\).
In this test \(L=(1,\lambda w_0)^*\mathcal P\), with
\(w_0+\bar w_0=1\) and \(\bar w_0\) the Rosati adjoint of \(w_0\).
Indeed these are the graph-line and adjunction conditions of the
datum. The graph identification can be made rigidified by its value
at zero. A rigidified isomorphism of cubical lines on an abelian
variety respects the cube structure by uniqueness.
\src{src/ReciprocityTheory/ReductiveGroups/ArithmeticGeometry/BirationalGeometry/ToroidalCompactification/TwoWeightEqOneAddRoot.lean}{82}{Graph-line condition};
\src{src/Formalism/LocalToGlobal/ProjectiveLattices/VAtVNeBot.lean}{53}{The integral model, eigendirections and theta generator}.

\begin{lemma}[The maximum of the canonical theta norm]\label{intjet:maximum}
On the analytification over the completed algebraic closure, the
canonical norm of \(s_0\) is at most one everywhere. Its maximum
is attained at the point of the skeleton belonging to the identity
component of the special fibre.
\end{lemma}
\begin{proof}
Use the valuation \(v=-\log|\cdot|\) with normalized uniformizer.
One may make a valued extension for this calculation. For the split
Raynaud uniformization let \(Y\) be the period lattice,
\(\Phi:Y\longrightarrow X\) its identification with the character
lattice under the principal polarization, and \(b\) the positive
definite monodromy form on \(Y_{\mathbb R}\).
For any rigidified ample line \(H\) giving \(\lambda\) take its
descent data on the Raynaud extension. Denote the constant of descent
at \(u\in Y\) by \(c_H(u)\) and the period pairing by \(t\), using
the convention
\[
 c_H(u+v)=c_H(u)c_H(v)t(u,\Phi v),\qquad
 -\log\|t(u,\Phi v)\|=b(u,v).
\]
Products in this identity use the biextension identifications of
the fibres. The Poincaré line pulled back by \(1\times\lambda\)
is \(m^*H\otimes p_1^*H^{-1}\otimes p_2^*H^{-1}\), with its
rigidification. Pulling back by the graph of \(w_0\) and taking
the quotient of the three descent constants therefore gives
\[
 c_L(u)=\frac{c_H(u+w_0u)}{c_H(u)c_H(w_0u)}
       =t(u,\Phi w_0u),\qquad
 -\log\|c_L(u)\|=b(u,w_0u)=\tfrac12 b(u,u).
\]
The last equality uses \(b(w_0u,v)=b(u,\bar w_0v)\), by functoriality
of the dual period pairing, and symmetry of \(b\).

Here are the precise standard uniformization results used to pass
from this computation to a metric. The balanced translate in
de Jong--Shokrieh, Proposition~8.3, has descent valuation
\(b(u,u)/2\). Translation by a point of tropical coordinate \(z\)
adds \(b(u,z)\) to the descent valuation (Lemma~8.1).
Thus the balancing point for the present triple has \(z=0\).
It acts trivially on the skeleton. The symmetry of both lines makes
its image a two-torsion point, since \(T_y^*L\simeq T_{-y}^*L\)
implies \(\lambda(2y)=0\). Their canonical metrics correspond under
translation up to a constant from rigidification (Lemma~5.1).
For the balanced translate Proposition~9.1 and the tropical
Riemann theta formula give on \(Y_{\mathbb R}/Y\)
\[
        -\log\|s_0(\nu)\|=\mathrm{constant}
                      +\tfrac12\min_{u\in Y} b(u+\nu,u+\nu).
\]
Its minimum occurs at \(\nu=0\bmod Y\). Corollary~10.4 gives
the same maximum of the norm over the whole analytification.
Proposition~10.3 at the identity component identifies the negative
log norm there with the order of the original section along that
component in its cubical model. This order is nonnegative by
integrality. This last application uses the original rigidification,
so no unspecified metric constant enters the inequality.
See de Jong and Shokrieh, \S\S6--11~\cite{deJongShokrieh}, and the
Raynaud theta descent in Foster, Rabinoff, Shokrieh and Soto,
\S\S3--4~\cite{fosterTheta}.
\end{proof}

\begin{lemma}[Integral traces on identity discs]\label{intjet:disc}
Let \(i:C_0\to C_1\) be an isogeny of abelian generic fibres induced
by an inclusion of full projective \(O_E\)-lattices tensored with
\(\mathcal A\). Write \(\mathfrak d R\) for the determinant ideal
of its differential on the integral Lie algebras. Let \(L_1\)
be a symmetric cubical line on \(C_1\) with its canonical metric
and integral cubical extension to the identity model. Suppose
\(u\) is a section of \(i^*L_1\) bounded by one everywhere
in the pullback metric. Write \(T=\operatorname{Tr}_i u\).
Take \(b\in C_1[h]\), where \(h\) is odd and prime to the residue
characteristic, choose \(c\) with \(hc=b\) and \(h^2c=0\),
and suppose the normalized translation at \(c\) exists on
\([h]^*L_1\). On the formal disc at zero the raw coefficient of
\(\widetilde T_c[h]^*T\) in an integral frame of \([h]^*L_1\)
is divisible by \(\mathfrak d R\).
\end{lemma}
\begin{proof}
Extend the ground field finitely to contain the points of the
relevant finite fibres. The homomorphism of completed identity
discs of the two models is finite flat. For details, each smooth
disc has ring \(R[[X_1,\ldots,X_g]]\). On special fibres the
kernel of \(i\) is finite since an integer multiplication factors
through \(i\) and acts as an isogeny of semiabelian varieties.
Consequently the images of target parameters and a uniformizer
generate an ideal primary to the maximal ideal of the source.
Lifting a basis of the finite quotient by this ideal and iterating
in the ideal gives finite generation over the complete target ring.
These \(g+1\) parameters form a regular sequence in the regular
local source. The local criterion for flatness, successively along
the parameters, gives a finite free module. This finite, finitely
presented flat map is a local complete intersection at the closed
point: its closed fibre is cut out by the regular sequence of \(g\)
specialized parameters in the smooth special-fibre disc. The complete
intersection locus is open, hence contains the whole local spectrum.
The finite module of relative differentials is computed by the
formal Jacobian. Indeed derivations into finite modules are
continuous for the maximal-ideal topology by the product rule and
these modules are complete. In integral translation-invariant
differential bases the Jacobian determinant has constant ideal
\(\mathfrak d R\). Thus the trace of any power series in the source
lies in \(\mathfrak d R[[Y]]\) by the different formula in
Theorem~\ref{intjet:trace}.

The rigid generic fibre of either completed local ring is its open
unit polydisc. The disc map just constructed is surjective there:
the finite free map has positive rank, and points lying above a
point reducing to zero reduce to zero in this local map. More
explicitly the primary-ideal identity gives, for some \(e\),
\(X_j^e\in(\pi,i_1(X),\ldots,i_g(X))\); thus every point of the
finite formal fibre above such a point has \(|X_j|<1\). Finite
base change, or evaluation first on closed subdiscs, gives the
analytic covering. It is étale in characteristic zero, and each
fibre is a coset of the subgroup \(K^\circ\subset\ker i\) lying in
the identity disc. The full preimage over that disc in the abelian
generic fibre is the disjoint union of its translates indexed by
\(\ker i/K^\circ\). Indeed choose a disc preimage of a given
point. All other preimages differ by elements of \(\ker i\), and
two translated discs meet precisely when their translates differ by
an element of \(K^\circ\).

An integral frame of a symmetric cubical line on the identity disc
has canonical norm one. To see the metric comparison, on the generic
fibre of the formal completion along the entire special fibre take
the model metric. This domain is covered by finitely many compact
affinoids and is stable under doubling. The rigidified cubical
isomorphism for \([2]\) implies that the logarithm \(F\) of the
ratio of canonical and model metrics satisfies \(F(2x)=4F(x)\).
It is bounded on this domain, so iteration gives \(F=0\).
An integral formal frame differing from a model frame by a power
series with unit constant term has the same norm on the identity
disc. The canonical metric is functorial for rigidified isogeny
pullbacks by the same doubling argument on the compact generic
fibre.

A finite-order lift \(U\) of translation at any torsion point
preserves the canonical metric on its symmetric line. Choose
\(m>1\) congruent to one modulo the order of the point. Translation
commutes with \([m]\); comparison of the two line isomorphisms
on the proper connected generic fibre differs by a scalar.
The continuous logarithmic distortion therefore satisfies
\(F([m]x)=m^2F(x)+\mathrm{constant}\). Iteration on the compact
space makes \(F\) constant, and the finite order of \(U\) makes it
zero. This also applies to the normalized lift at \(c\).

In the identity frame on the target, view the translated pullback
as the trace above \(y\mapsto b+[h]y\), using that lift to identify
the line. Multiplication by \(h\) is an isomorphism on either
identity disc, by the integral formal inverse function theorem.
For each class in \(i^{-1}(b)/K^\circ\), choose a representative
\(x\). The preimages on its sheet are parametrized by
\(x+[h]z\), with \(i(z)=y\). On this disc the coefficient of
\(u(x+[h]z)\) in the pulled-back translated frame is analytic of
absolute value at most one, by the metric calculation. Every such
function on the open unit polydisc has an integral power series:
apply the Gauss norm bound on closed polydiscs of radii tending to
one, after a valued extension if needed. The coefficients already
belong to the field of definition. The sum over this sheet is the
trace of that power series under the finite free formal map \(i\).
On geometric generic fibres this is the usual sum over the sheet,
which identifies the two analytic traces. Each is divisible by
\(\mathfrak d R\), so their finite sum has the same divisibility.
\end{proof}

\begin{theorem}[Finite denominators at semiabelian tests]\label{intjet:semi}
At every integral test with the stated semiabelian, cubical and
Hodge data, \(D_M^2\) clears the rescaled finite theta sum. The
valuation inequalities of Proposition~\ref{intjet:standard} at
locally standard lattices and of Theorem~\ref{intjet:trace} away
from \(M\) hold also at semiabelian reduction.
\end{theorem}
\begin{proof}
At a locally standard lattice use \(N_0\) and the generic-fibre
identity of Proposition~\ref{intjet:standard}. The norm of the
generator is at most one everywhere by Lemma~\ref{intjet:maximum}.
Its normalized pullback by \([N_0]\) and its two finite-order
translates have the same bound, by functoriality and the isometry
proved in Lemma~\ref{intjet:disc}. In the integral formal frame
at zero their raw germs therefore have integral coefficients by
the polydisc argument. Lemma~\ref{intjet:der} applies on a smooth
semiabelian formal group verbatim: it uses the formal group law,
the integral tangent directions and the cubical adjunction.
This proves the individual inequality, also when \(\ell\mid M\).

For \(\ell\nmid M\) use the notations \(a_0,i,j,a\) in
Theorem~\ref{intjet:trace}. The tensor models are direct summands
of products of the smooth semiabelian scheme. An idempotent on a
smooth group gives a product decomposition by its image and the
image of its complement: these are the closed kernels of the
opposite projections, and addition and the pair of projections
are mutual inverses. Each factor is smooth by the lifting
criterion and has semiabelian fibres as a direct factor of such a
group. The graph-line construction on these models uses only the
integral cubical biextension and its \(O_E\)-adjunction. On the
generic fibre the identities of pullback lines and normalized
translations give, after summing the first translation,
\(a^*\operatorname{Tr}_i(j^*s_0^{\boxtimes n})\), as before.
The section inside this trace has norm at most one everywhere,
since the rigidified line identifications are isometries.

The final point \(t\) has odd order \(h\) prime to \(\ell\).
Take \(t'\) with \(ht'=t\), and set \(c=a(t')\), \(b=a(t)\).
The normalized lift at \(c\) on \([h]^*\mathcal L_1\) exists:
\(h^2c=0\) and \([h]^*\mathcal L_1=\mathcal L_1^{h^2}\).
Set \(T=\operatorname{Tr}_i(j^*s_0^{\boxtimes n})\).
The sections
\[
       [h]^*\widetilde T_t\,a^*T,\qquad
           a^*\widetilde T_c\,[h]^*T
\]
agree up to a root of unity. Both arise from the same translation
\(t'\) on \([h]^*a^*\mathcal L_1\), by pulling back the respective
lifts and using \([h]T_{t'}=T_t[h]\) and \(aT_{t'}=T_ca\).
These two lifts have finite order: a power divisible by the point
order and the order of the original lift is the identity. Their
ratio is a scalar on the proper generic fibre, hence is a root
of unity. Lemma~\ref{intjet:disc} gives divisibility by
\(\mathfrak b R\) for the section on the right before applying
\(a^*\), since the determinant on integral Lie algebras has
exactly the same ideal computed in Theorem~\ref{intjet:trace}.
Multiplication by \(s\in\mathfrak b^{-1}\) gives an integral raw
germ. Apply the integral polynomial and harmonic chain rule for
\(a=U(NV)\) in that theorem. The extra pullback \([h]\) on the
left multiplies a homogeneous corrected jet by \(h^d\), a unit.
The factor \(N^d\) and the loss of at most
\(\lfloor d/2\rfloor v(\Delta)\) follow as claimed. The index
cancellation and the level factor \(D_M=C_M(M\Delta)^d\) finish
the original finite-place test.
\end{proof}
